\documentclass[11pt]{article}

\usepackage[margin=1in]{geometry}
\usepackage{array}
\usepackage{booktabs}
\usepackage{longtable}
\usepackage{hyperref}
\usepackage{xcolor}

\title{HAOS-IIP Scale-Bridge Legitimacy Audit:\\A Phase 66-Compatible Companion to the Continuum Physics Scaling Roadmap}
\author{Tomislav Rupi\'c \\ Independent Researcher}
\date{May 2026}

\begin{document}

\maketitle

\begin{abstract}
This companion note freezes the scale-bridge legitimacy audit for HAOS-IIP as a Phase 66-compatible hostile-audit artifact. It is not a new phase, not Phase 67, not a new mechanism, and not a physical correspondence claim. Its purpose is to strengthen the existing Continuum Physics Scaling Roadmap by formalizing how HAOS-IIP evaluates refinement-ordered scaling diagnostics toward continuum feasibility. The central rule is simple: refinement evidence is not continuum proof. The audit defines ledger requirements for scale-facing claims, matched-control scaling comparisons, proto-geometry surrogate language rules, limited metric-like tests, coarse-graining recovery checks, status-box requirements, corrected public language, and explicit failure gates. It also records a v1.1 focused lift-cycle scaffold: initial convergence diagnostics in progress, matched-control and coarse-graining templates, metric-surrogate templates, and a bounded predictive bridge template. No numerical convergence result is introduced here. The document raises the burden of proof before any stronger continuum-facing statement can be made.
\end{abstract}

\section{Status}

This paper is a numbered PDF companion to:
\begin{itemize}
\item \texttt{docs/notes/foundations/HAOS\_IIP\_Scale\_Bridge\_Legitimacy\_Audit\_v1.md}
\item \texttt{docs/notes/foundations/HAOS\_IIP\_Continuum\_Physics\_Scaling\_Roadmap\_v1.md}
\item \texttt{docs/notes/foundations/HAOS\_IIP\_Phase\_66\_Translation\_and\_Hostile\_Audit\_Layer\_v1.md}
\end{itemize}

It is:
\begin{itemize}
\item a Continuum Physics Scaling Roadmap companion artifact;
\item a Phase 66-compatible hostile-audit document;
\item a scale-bridge legitimacy checklist;
\item a public-language correction layer for continuum-facing claims.
\end{itemize}

It is not:
\begin{itemize}
\item a new phase sequence;
\item Phase 67;
\item a new operator mechanism;
\item a proof of a continuum limit;
\item a derivation of spacetime, known physics, physical geometry, or physical correspondence.
\end{itemize}

\section{Core Rule}

\begin{center}
\fbox{\parbox{0.86\linewidth}{
\centering
\textbf{Refinement evidence is not continuum proof.}
}}
\end{center}

The scale bridge becomes more legitimate only when it survives stricter bookkeeping, matched controls, refinement, and compression. The audit does not make the bridge more dramatic. It makes it harder to overstate.

\section{Current Scale-Bridge Problem}

HAOS-IIP has strong internal frozen-regime diagnostics, scalar-carrier geometry closure on one tested kernel-graph family, bounded proto-continuum evidence, and a reproducible late-phase propagation, ordering, causal-closure, and distance-surrogate chain. That is meaningful. It is not yet a proven continuum limit.

The current weakness is the scale bridge. The program can show refinement-ordered behavior and metric-like surrogate stabilization inside tested operator regimes, but it has not shown that these structures converge uniquely, universally, or physically toward a continuum geometry or continuum field theory.

The audit therefore asks:
\begin{quote}
Is HAOS-IIP detecting structured refinement behavior, or generic graph behavior with disciplined language?
\end{quote}

The honest current status is:
\begin{itemize}
\item internal diagnostics are strong inside declared frozen regimes;
\item CP4 has a bounded shared-family scalar kernel-graph geometry closure;
\item proto-geometric surrogate behavior exists at the tested level;
\item CP5 universality is still open;
\item continuum ontology is not established.
\end{itemize}

\section{Relation to the CP Ladder}

This audit attaches to the existing CP1--CP6 ladder. It does not introduce a separate sequence.

\begin{longtable}{p{0.18\linewidth}p{0.34\linewidth}p{0.38\linewidth}}
\toprule
CP gate & Roadmap role & Scale-bridge audit question \\
\midrule
CP1 & Scalar continuum contract & Does the scalar operator converge beyond one clean substrate family and one normalization? \\
CP2 & Vector continuum contract & Is apparent vector behavior a true low coexact branch or inherited/topological structure? \\
CP3 & Effective-equation contract & Do coarse laws survive refinement without retuning and outperform controls? \\
CP4 & Geometry closure contract & Does one metric-like surrogate organize the same tested branch family without importing geometry by hand? \\
CP5 & Universality contract & Does the effective statement survive nontrivial changes of construction? \\
CP6 & Bounded continuum-physics statement & Can the result be stated as bounded continuum feasibility without ontology creep? \\
\bottomrule
\end{longtable}

CP1, CP4, and CP5 carry the heaviest scale-bridge load. CP1 controls the scalar continuum foothold. CP4 controls whether geometry-like language is earned inside one shared operator/substrate family. CP5 controls whether the bridge is a one-family artifact or a robust continuum-feasibility claim. If CP5 fails, CP6 must shrink.

\section{Refinement Scaling Ledger}

Every continuum-facing scale claim should create or reference a public refinement scaling ledger. A clean-looking trend without its failed controls is not scale-bridge evidence.

\begin{longtable}{p{0.27\linewidth}p{0.63\linewidth}}
\toprule
Field & Meaning \\
\midrule
\texttt{claim\_id} & Stable identifier for the scale-bridge claim \\
\texttt{operator\_family} & \texttt{L0}, \texttt{L1}, \texttt{D\_H}, block cochain Laplacian, scalar carrier, or declared operator \\
\texttt{substrate\_graph\_family} & Cubic grid, perturbed grid, random geometric graph, kernel graph, periodic box, or declared family \\
\texttt{refinement\_level} & \texttt{n\_side}, \texttt{h}, sample count, or hierarchy index \\
\texttt{kernel\_width\_scaling\_rule} & Fixed width, \(\epsilon_k \sim h\), \(\epsilon_k \sim h^2\), adaptive width, or other rule \\
\texttt{normalization} & Discrete second-moment normalization, trace normalization, mass normalization, or declared alternative \\
\texttt{measured\_observable} & Eigenvalue, trace ratio, Green response, shell slope, metric surrogate, current closure, etc. \\
\texttt{matched\_control} & Explicit control family used at the same level \\
\texttt{error\_trend} & Refinement-ordered error, residual, drift, or fit gap \\
\texttt{dimensionless\_ratio} & Unitless ratio used to compare across refinement \\
\texttt{recoverability\_index} & Retention or reconstruction score under perturbation \\
\texttt{branch\_identity\_persistence} & Overlap, principal angle, label stability, or spectral address persistence \\
\texttt{metric\_surrogate\_stability} & Metric-like stability score, bounded violation rate, or ordering stability \\
\texttt{status} & \texttt{PASS}, \texttt{OPEN}, or \texttt{FAIL} \\
\texttt{claim\_boundary} & What the row explicitly does not claim \\
\bottomrule
\end{longtable}

\section{Matched-Control Scaling Comparison}

Every scale-bridge claim must compare the branch result against controls chosen to be close enough to be dangerous:
\begin{itemize}
\item random graph control;
\item shuffled-weight control;
\item matched-degree control;
\item non-admissible operator control;
\item perturbed-kernel control;
\item seed-randomized ensemble.
\end{itemize}

Control comparison should happen at the same level as the claim. If the claim concerns eigenvalue scaling, compare eigenvalue scaling. If it concerns metric-like behavior, compare metric-like behavior. If it concerns coarse-law closure, compare equation-fit residuals and coefficient drift.

\section{Proto-Geometry Surrogate Language}

Allowed standard-language terms:
\begin{itemize}
\item metric surrogate;
\item distance proxy;
\item curvature proxy;
\item spectral separation;
\item subspace overlap decay;
\item Green-response distance;
\item heat-kernel behavior;
\item shell-arrival slope;
\item operator-response distance;
\item refinement-ordered scaling diagnostic.
\end{itemize}

Forbidden stronger terms unless later proven by a stronger contract:
\begin{itemize}
\item emergent spacetime;
\item derived physical geometry;
\item physical metric;
\item real curvature;
\item continuum manifold;
\item field-theory derivation.
\end{itemize}

Correct public wording:
\begin{quote}
metric-like surrogate behavior stabilizes under refinement inside the tested operator regime
\end{quote}

Incorrect public wording:
\begin{quote}
proto-geometry emerges
\end{quote}

unless the document has already specified the surrogate, controls, gates, and non-claims.

\section{Limited Metric-Like Tests}

For each metric surrogate or distance proxy, the document must test or explicitly defer:
\begin{itemize}
\item non-negativity;
\item identity behavior within numerical tolerance;
\item symmetry, when applicable;
\item triangle-inequality behavior or bounded violation;
\item refinement stability;
\item perturbation stability;
\item matched-control separation.
\end{itemize}

Passing these tests means metric-like surrogate behavior stabilizes inside the tested operator regime. It does not mean HAOS-IIP has derived spacetime or a physical metric.

\section{Coarse-Graining Recovery}

Refinement alone is not enough. The structure must also survive compression. For refinement levels \(h_i<h_j\), with projection or coarse-grain maps
\[
\Pi_{h_i\rightarrow h_j},
\]
the audit requires testing whether refined structure projects back to coarser levels while preserving:
\begin{itemize}
\item branch identity;
\item recoverability ranking;
\item spectral address persistence;
\item metric-surrogate ordering;
\item trace-ratio stability.
\end{itemize}

Core question:
\begin{quote}
Does the structure survive both refinement and compression?
\end{quote}

Failure under compression means the apparent refinement trend may be a high-resolution artifact.

\section{Scale-Bridge Status Box}

Every future continuum-facing document should include this status box:

\begin{quote}
\textbf{Scale-Bridge Status Box}

What is refined: \(\ldots\)

What stabilizes: \(\ldots\)

What fails: \(\ldots\)

What matched control was used: \(\ldots\)

What continuum-feasibility evidence exists: \(\ldots\)

What is not claimed: \(\ldots\)
\end{quote}

The final field is mandatory. This does not claim a proven continuum limit, derivation of spacetime, derivation of known physics, physical correspondence, uniqueness of the limit, or replacement of established models.

\section{Corrected Language Rules}

\begin{longtable}{p{0.39\linewidth}p{0.50\linewidth}}
\toprule
Replace & Use instead \\
\midrule
continuum bridge & refinement-ordered scaling diagnostics toward continuum feasibility \\
proto-geometry emerges & metric-like surrogate behavior stabilizes under refinement inside the tested operator regime \\
physical observable & bridge observable candidate \\
derived metric field & operator-derived metric surrogate \\
\bottomrule
\end{longtable}

Short phrases are allowed in titles and file names, but the first public explanation must use the corrected language.

\section{Failure Gates}

A scale-bridge claim fails if:
\begin{itemize}
\item measured structure does not stabilize under refinement;
\item stabilization is no better than matched controls;
\item metric-like behavior depends on hand-picked parameters;
\item branch identity changes unpredictably across refinement;
\item coarse-graining destroys the claimed structure;
\item error trends do not improve with refinement;
\item null models reproduce the same effect;
\item interpretation requires importing geometry by hand.
\end{itemize}

When a failure gate triggers, the claim must shrink to the last surviving lower gate. A failed CP5 universality claim may still leave a valid CP4 one-family geometry closure. A failed metric surrogate may still leave valid operator convergence.

\section{v1.1 Focused Lift-Cycle Templates}

Status:
\begin{itemize}
\item focused lift cycle active;
\item initial convergence diagnostics in progress;
\item no new numerical result recorded in this audit yet;
\item target: complete the first minimal convergence and matched-control diagnostic before Phase XIX freeze.
\end{itemize}

\subsection{Minimal convergence table template}

\begin{center}
\small
\begin{tabular}{p{0.08\linewidth}p{0.15\linewidth}p{0.15\linewidth}p{0.14\linewidth}p{0.14\linewidth}p{0.12\linewidth}p{0.10\linewidth}}
\toprule
Level & Operator & Observable & HAOS value & Control value & Recovery & Status \\
\midrule
TBD & TBD & spectral gap / trace ratio & TBD & TBD & TBD & OPEN \\
\bottomrule
\end{tabular}
\end{center}

Minimum acceptance rule: at least three deterministic refinement levels, one matched control with the same observable and normalization, one coarse-graining recovery score, explicit PASS/OPEN/FAIL, and explicit claim boundary in every row.

\subsection{Matched-control failure example template}

\begin{center}
\small
\begin{tabular}{p{0.20\linewidth}p{0.20\linewidth}p{0.22\linewidth}p{0.23\linewidth}}
\toprule
Control family & What was matched & What failed & Interpretation \\
\midrule
TBD & degree / weights / seed / kernel / size & scaling trend / branch identity / recovery & control does not reproduce the claimed refinement-ordered structure \\
\bottomrule
\end{tabular}
\end{center}

The control failure must state what was held fixed, what diverged, and whether the difference survives normalization.

\subsection{Scale-Bridge Status and Limits paragraph template}

\begin{quote}
This artifact tests refinement-ordered scaling diagnostics toward continuum feasibility inside the declared HAOS-IIP operator regime. What is refined is [operator/substrate/level]. What stabilizes is [observable/signature/surrogate]. What fails or remains open is [failure/open boundary]. The matched control is [control family]. The continuum-feasibility evidence is [bounded evidence]. This does not claim a proven continuum limit, derivation of spacetime, derivation of known physics, physical correspondence, uniqueness of the limit, or replacement of established models.
\end{quote}

\subsection{Bounded predictive bridge result template}

\begin{center}
\small
\begin{tabular}{p{0.15\linewidth}p{0.20\linewidth}p{0.13\linewidth}p{0.18\linewidth}p{0.13\linewidth}p{0.13\linewidth}}
\toprule
Bridge line & Predeclared target & Threshold & Data boundary & Result & Status \\
\midrule
ferron / magnon / biology & spectral feature / recovery threshold / target window & TBD & raw-data status & TBD & OPEN \\
\bottomrule
\end{tabular}
\end{center}

Minimum acceptance rule: target declared before full audit, raw-data boundary stated, failures recorded next to successes, and no language stronger than bridge observable candidate.

\section{Final Boundary}

Scale-bridge evidence is useful.

Scale-bridge evidence is not continuum proof.

The correct posture is not pessimism. It is recoverable coherence under stronger interaction with controls, refinement, and compression.

\end{document}
